\begin{abstract}

Khóa luận này trình bày quá trình \textbf{tìm hiểu, thiết kế và hiện thực một
nền tảng blockchain có cấp phép (permissioned blockchain) định hướng cho doanh
nghiệp}, được xây dựng hoàn toàn từ đầu (from-scratch) và lấy cảm hứng từ kiến
trúc của Hyperledger Fabric. Khác với phần lớn các đồ án sử dụng lại một nền
tảng blockchain có sẵn, mục tiêu cốt lõi của chúng tôi là \textit{làm chủ công
nghệ lõi}: tự cài đặt từng thành phần để hiểu thấu đáo cách một hệ thống sổ cái
phân tán hoạt động, từ tầng mạng ngang hàng, thuật toán đồng thuận, máy ảo hợp
đồng thông minh, cho đến tầng lưu trữ bất biến và world state.

Hệ thống được tổ chức theo triết lý tách biệt ba trách nhiệm
\textbf{Execute--Order--Validate} của Hyperledger Fabric, gồm bốn thành phần độc
lập giao tiếp qua mạng ngang hàng \texttt{libp2p}: (1)~\textbf{Core Service}
đóng vai trò cổng vào và tầng tính toán, thực thi hợp đồng thông minh biên dịch
sang WebAssembly trong môi trường cô lập \texttt{wazero} và điều phối quá trình
xác nhận (endorsement); (2)~\textbf{Ordering Service} --- trái tim đồng thuận
của hệ thống --- \textit{tự cài đặt thuật toán Raft từ đầu} với một biến thể bầu
lãnh đạo theo độ ưu tiên (priority-based) thay cho cơ chế timeout ngẫu nhiên;
(3)~\textbf{Committing Peer} kiểm tra tính hợp lệ về mật mã của block và ghi
vĩnh viễn vào sổ cái append-only kèm world state UTXO trên LevelDB; và
(4)~\textbf{Blockchain Explorer} là giao diện web React hiển thị dữ liệu theo
thời gian thực qua Server-Sent Events. Một cơ sở dữ liệu PostgreSQL đóng vai trò
bản sao (mirror) bất đồng bộ phục vụ tra cứu và đo lường, hoàn toàn nằm ngoài
luồng đồng thuận.

Khóa luận đặc biệt đi sâu vào tầng đồng thuận: mô hình hóa các trạng thái của
node, quy trình bầu lãnh đạo xác định, cơ chế nhân bản log và cắt block theo cơ
chế lai (hybrid) giữa kích thước lô và thời gian, quá trình giao block (deliver)
theo mô hình đăng-ký--phát-tỏa, và cơ chế đồng bộ (sync) bốn pha dựa trên đồng
thuận \texttt{(commitIndex, commitHash)}. Chúng tôi cũng trình bày một loạt tối
ưu hiệu năng trên đường nóng (hot path) giúp hệ thống tiệm cận mục tiêu
\textbf{5.000 giao dịch mỗi giây}, cùng phương pháp đo lường ba lớp (k6, submit,
commit) để phân tích throughput, độ trễ end-to-end và hiện tượng backlog dưới
tải bão hòa.

\vspace{0.4cm}
\noindent\textbf{Từ khóa:} blockchain có cấp phép, Hyperledger Fabric, đồng
thuận Raft, hợp đồng thông minh WebAssembly, mạng ngang hàng libp2p, mô hình
Execute--Order--Validate, sổ cái phân tán, throughput, độ trễ.

\end{abstract}
